\documentclass[10pt,a4paper]{amsart}
\usepackage[margin=19mm]{geometry}
\usepackage{amsmath,amssymb,mathtools}
\usepackage[scr=rsfs,cal=euler]{mathalfa}
\usepackage{xfrac}
\usepackage[unicode,hidelinks,bookmarks=false]{hyperref}
\newcommand{\ie}{i.e.\ }
\newcommand{\cf}{cf.\ }
\newcommand{\resp}{respectively\ }
\newcommand{\Subsection}{\subsection}
\providecommand{\nobreakdash}{\nobreak\hbox{-}}
\setlength{\emergencystretch}{2em}
\multlinegap0pt
\usepackage{bm}
\usepackage{bbm}
\usepackage{dsfont}
\usepackage{xspace}
\usepackage{cancel}
\usepackage{mathtools}
\usepackage{amsthm}
\makeatletter
\@ifundefined{theorem}{\newtheorem{theorem}{Theorem}}{}
\@ifundefined{lemma}{\newtheorem{lemma}[theorem]{Lemma}}{}
\makeatother
\newcommand{\C}{\mathbb{C}}
\title[On stability of outliers from the circular law]{On stability of outliers from the circular law}
\author{Guillaume Dubach, Aniss Fares, Lilou Thomas}
\date{Source: arXiv:2606.16609v1 (CC BY 4.0)}
\begin{document}
\maketitle

\noindent\emph{Source and licence.} On stability of outliers from the circular law. Version arXiv:2606.16609v1 (CC BY 4.0), distributed under the Creative Commons Attribution 4.0 International licence, as recorded in the arXiv licence metadata for this exact version and shown on the abstract page https://arxiv.org/abs/2606.16609v1. Full rights notice: licenses/arxiv-2606.16609-CC-BY-4.0.txt.\par\smallskip

\noindent\textit{Editorial scope.} Counted unit: the Lemma whose source label is \texttt{None} in the article (source0.tex lines 461--467, source label \texttt{None}), delivered together with the article's own complete original proof of it (source0.tex lines 468--519, one \texttt{proof} environment of 52 lines). No mathematics is written, repaired or re-derived: statement and proof are the article's own text line for line. Required context retained uncounted: none. Every definition, display and label of those lines is retained unchanged. Omitted from this package and not redistributed: 1--460, 520--1105. Nothing inside a retained excerpt is omitted or altered: every retained body is delivered exactly as the article's lines are written, so no omission receipt of any kind accompanies this package. Citations: the retained excerpts carry no citation command at all, so no bibliography entry of the article is transcribed and no reference is redistributed. Reference adaptation: none was needed and none was made; every reference inside the delivered unit resolves inside it, so no unresolved reference or citation is produced. Printed slips kept literal: no printed word, symbol, hypothesis or number of the counted statement or of its proof was altered. Layout and class: the article's own class is \texttt{amsart} with packages including geometry, graphicx, tikz, amsmath, mathrsfs, physics, amssymb, amsrefs, bm, dsfont, hyperref; the wrapper is the lane's pinned \texttt{amsart} editorial wrapper at 10pt on A4, which restates the theorem-like environments the retained text needs and adds the article's own macro definitions that the retained text needs (\texttt{\textbackslash C}), with extra wrapper packages (bm, bbm, dsfont, xspace, cancel, mathtools). The delivered numbering is the unit's own. Assets: no retained span contains a figure environment, a graphics-inclusion command, a PDF-page inclusion, a table or a TikZ picture; no image file is copied into this package, the package declares no asset dependency and no image file is redistributed with it. Licence: version arXiv:2606.16609v1 of this article is distributed under the Creative Commons Attribution 4.0 International licence, as recorded in the arXiv licence metadata for this exact version and shown on the abstract page https://arxiv.org/abs/2606.16609v1. A full rights notice is delivered at licenses/arxiv-2606.16609-CC-BY-4.0.txt. Characters: every retained excerpt is pure ASCII, so the delivered engine, XeLaTeX with its default Latin Modern text font, reports no missing character.
\begin{lemma}
   Set $Q_N(z):=1+b_N^*\mathbf{R}_{N}(z)\mathbf{R}_{N}(z)^*b_N$ for $z \notin Sp (D_N)$ and denote by $\mu_1,\ldots,\mu_{N-1}$ the eigenvalues of $D_N$. Then 
   \begin{equation*}
       Q_N(z) \overset{}{=} \prod_{k=1}^{N-1}\left(1+\frac{\lvert X_k\lvert^2}{N\lvert z- \mu_k \lvert^2}\right)
   \end{equation*}
   where $X_1,\ldots,X_{N-1}$ are i.i.d. standard complex Gaussian random variables, independent of $\mu_1,\ldots,\mu_{N-1}.$
\end{lemma}

\begin{proof}
    Let $D_N=UTU^*$ be a Schur decomposition of $D_N$, where $T$ is upper triangular and has diagonal entries $\mu_1,\ldots, \mu_{N-1}$. Since $b_N$ is independent of $D_N$ and has a unitary invariant Gaussian $\beta:=U^*b$ has law $\mathcal{N}_{C}(0,N^{-1}I)$ and is independent of $T$ and moreover 
    \begin{equation*}
        b_N^*\mathbf{R}_{N}(z)\mathbf{R}_{N}(z)^*b_N=\beta ^*(z-T)^{-1}(z-T)^{-*}\beta.
    \end{equation*}
    Thus it is enough to prove the identity for $T$ and $\beta$. \\
    Let $T_k $ be the $k\times k$ upper-left principal block of $T$, and let $\beta_k$ be the vector made of the first $k$ coordinates of $\beta.$ Define $R_k(z):=(z-T_k)^{-1}$ and $Q_k(z):=1+\beta_k^* R_k(z)R_k(z)^*\beta_k.$ Then $Q_0(z)=1$ and $Q_{N-1}(z)$ is the quantity of interest. Write 
    \begin{equation*}
        T_{k+1}= \begin{pmatrix}
        T_k&\tau_k\\
        0&\mu_{k+1}
    \end{pmatrix}, \ \beta_{k+1}=\begin{pmatrix}
        \beta_k\\
        \eta_{k+1}
    \end{pmatrix}.
    \end{equation*}
    Conditionally on the spectrum, the vectors $\tau_k$ have independent complex Gaussian entries of variance $1/N$ and are independent of the previous columns. Also, the variables $\eta_k$ are independent complex Gaussians of variance $1/N$, independent of $T$. Set $y_k:=R_k(z)^*\beta_k.$ Then $Q_k(z)=1+ \lvert \lvert y_k \lvert \lvert ^2.$ Write 
    \begin{equation*}
        y_{k+1}=\begin{pmatrix}
        y_k\\
        \gamma_{k+1}
    \end{pmatrix}.
    \end{equation*}
    Since $(z-T_{k+1})^*y_{k+1}=\beta_{k+1}$ and 
    \begin{equation*}
        (z-T_{k+1})^*= \begin{pmatrix}
        (z-T_k)^*&0\\
        -\tau_k^*& \overline{z}-\overline{\mu_{k+1}}
    \end{pmatrix}
    \end{equation*}
    we obtain that 
    \begin{equation*}
        \gamma_{k+1}=\frac{\eta_{k+1}+\tau_k^*y_k}{\overline{z}-\overline{\mu_{k+1}}}.
    \end{equation*}
    Therefore, 
    \begin{equation*}
        Q_{k+1}(z)=1+\lvert \lvert y_{k+1}\lvert \lvert^2=Q_k(z)+ \frac{\lvert \eta_{k+1}+\tau_k^*y_k\lvert^2}{\lvert z-\mu_{k+1}\lvert^2}
    \end{equation*}
    Now define $X_{k+1}:=\sqrt{\frac{N}{Q_k(z)}}(\eta_{k+1}+\tau_k^*y_k).$ Then 
    \begin{equation*}
        Q_{k+1}(z)=Q_k(z)\left(1+ \frac{\lvert X_{k+1}\lvert^2}{N\lvert z-\mu_{k+1}\lvert^2}\right).
    \end{equation*}
    We now identify the law of the $X_k$. Let $\mathcal{F}_k:=\sigma(\mu_1,\ldots,\mu_{N-1},T_k,\beta_k).$ Then $y_k$ and $Q_k(z)$ are $\mathcal{F}_k$-measurable. Moreover, $\eta_{k+1}$ and $\tau_k$ are independent of $\mathcal{F}_k$ and 
    \begin{equation*}
        \eta_{k+1}+\tau_k^*y_k \mid \mathcal{F}_k \sim \mathcal{N}_{\C}\left(0,\frac{1+\lvert \lvert y_k \lvert \lvert^2}{1}\right)=\mathcal{N}_{\C}\left(0,\frac{Q_k(z)}{N}\right).
    \end{equation*}
    Hence $X_{k+1}\mid \mathcal{F}_k \sim \mathcal{N}_{C}(0,1).$ The conditional law does not depend on $\mathcal{F}_k$, therefore $X_{k+1}$ is independent of $\mathcal{F}_k$. In particular, $X_{k+1}$ is independent of the spectrum and of $X_1,\ldots,X_k$. By induction $X_1,\ldots,X_{N-1}$ are i.i.d. standard complex Gaussians independent of $\mu_1,\ldots,\mu_{N-1}$. Iterating the recurrence from $Q_0(z)=1$, we obtain 
    \begin{equation*}
        Q_{N-1}(z)=\prod_{k=1}^{N-1}\left(1+\frac{\lvert X_k \lvert ^2}{N\lvert z-\mu_k \lvert^2}\right)
    \end{equation*}
    which concludes the proof.
\end{proof}

\end{document}
